% !TEX program = xelatex
% Шаблон научной статьи НИТУ «МИСиС»
% Заполните блоки ниже и скомпилируйте: xelatex article.tex

\documentclass{misisarticle}

% --- Классификационные коды (оставьте пустым {}, если не требуются) ---
\grnti{28.17.19}
\asjc{2611, 3309}
\oecd{1.01, 1.02}

% --- Название ---
\title{Оптимизация защиты тарифных классов на регулярном авиарейсе методом имитационного моделирования}

% --- Авторы: \misisauthor{ФИО}{аффилиация}{город, страна}{email} ---
\misisauthor{З.~А.~Петроченко}{кафедра инженерной кибернетики НИТУ «МИСиС»}{Москва, Россия}{m2208598@edu.misis.ru}
\misisauthor{А.~С.~Кожаринов}{кафедра инженерной кибернетики НИТУ «МИСиС»}{Москва, Россия}{kozharinov@misis.ru}

% --- Аннотация и ключевые слова ---
\abstract{%
  Работа посвящена проблеме максимизации прибыли авиакомпании за счёт механизма защиты тарифных классов. Для реализации механизма предлагается оценка параметров кривой спроса, восстановленного по накопленным бронированиям с учётом ограничений на продажи. Лимиты защиты определяются методом имитационного моделирования Монте-Карло при стохастическом спросе. Задача решается на исторических данных одной из авиакомпаний. Качество модели оценивается сравнением выручки, полученной за рассматриваемый период по историческим данным, и выручки, рассчитанной вновь для тех же рейсов.
}

\keywords{%
  управление ресурсами авиакомпании, имитационное моделирование, лимиты бронирования, защита тарифных классов, метод Монте-Карло, стохастический спрос
}

\begin{document}

\maketitle

% ============================================================
\section{Введение}

Управление доходами (revenue management) авиакомпании~--- это совокупность методов, позволяющих продавать каждое место на рейсе нужному клиенту, в нужное время и по подходящей цене~\cite{talluri2004}. Один и тот же салон самолёта продаётся в нескольких тарифных классах, цена которых различается в несколько раз. Пассажиры, готовые платить больше, как правило, покупают билеты поздно, а чувствительные к цене~--- заранее. Если открыть дешёвые классы без ограничений, места будут распроданы раньше, чем появятся покупатели дорогих билетов, и авиакомпания недополучит выручку. Если же, напротив, закрыть дешёвые классы слишком рано, часть мест останется незаполненной.

Для решения этой проблемы применяется механизм \textit{защиты тарифных классов}: для каждого более дорогого класса резервируется часть мест, недоступная для продажи в более дешёвых классах. Размер этой части определяется лимитами бронирования. Качество лимитов зависит от того, насколько точно оценён спрос на каждый класс и насколько корректно учтена его случайная природа.

Практическая сложность состоит в том, что наблюдаемые по инвентарным данным бронирования не равны истинному спросу. Как только класс закрыт для продажи, дальнейший спрос на него в данных не отражается (цензурирование спроса). Поэтому оценка параметров спроса по сырым данным приводит к его занижению и к смещённым лимитам.

Целью работы является разработка и проверка на реальных данных метода выбора лимитов защиты тарифных классов на одном рейсе, основанного на восстановлении спроса по кривым накопления бронирований и имитационном моделировании.

Статья имеет следующую структуру. В разделе~II рассматриваются известные работы по теме. В разделе~III формулируется задача. В разделе~IV описываются методы: модель спроса, имитационная модель продаж, поиск оптимальных лимитов и способ валидации.

% ============================================================
\section{Анализ работ}

Классическая постановка задачи защиты мест для двух тарифных классов предложена в работе К.~Литтлвуда~\cite{littlewood1972}. Согласно его правилу, место следует защитить для дорогого класса, пока вероятность того, что спрос на него превысит текущий уровень защиты, не меньше отношения цен дешёвого и дорогого классов. Для нескольких вложенных классов используются эвристики семейства EMSR, предложенные П.~Белобой~\cite{belobaba1987}. Подробное изложение аналитических методов и их свойств содержится в монографии~\cite{talluri2004}. Эти методы опираются на предположение об известном и, как правило, нормальном распределении спроса по классам, и параметры этих распределений должны быть оценены независимо.

Имитационный подход позволяет отказаться от ряда допущений аналитических методов и проверить политику лимитов на модели процесса продаж. В работе Г.~М.~Фридмана и В.~М.~Баранова~\cite{fridman2025} построена имитационная модель процесса продаж авиабилетов для продуктов одного рейса. Вопросы моделирования процесса управления доходами на уровне рейсов рассмотрены Е.~В.~Носовой~\cite{nosova}, а описание компьютерного симулятора для системы управления доходами приведено в работе О.~В.~Лавренюк~\cite{lavrenyuk}. Задача управления сверхлимитными продажами с использованием математического моделирования исследована К.~А.~Мозговой~\cite{mozgovaya2013}.

Анализ перечисленных работ позволяет выделить три нерешённых вопроса, на которые направлено данное исследование.

\begin{enumerate}[label=\arabic*)]
  \item Методы либо описаны теоретически, либо проверены на модельных или учебных данных, которые нельзя напрямую сопоставить с работой действующей авиакомпании. Отсутствует сравнение фактически применённой на рейсе политики лимитов с оптимизированной по имитации политикой на одних и тех же исторических рейсах и по единой метрике~--- ожидаемой выручке, с доверительными интервалами по результатам Монте-Карло моделирования.
  \item Спрос, как правило, задаётся или предполагается заранее, а не восстанавливается по реальным кривым накопления бронирований.
  \item Валидация модели чаще выполняется на симулированных данных, а не на фактических показателях выручки того же исторического периода.
\end{enumerate}

% ============================================================
\section{Постановка задачи}

Рассматривается один регулярный рейс с вместимостью $C$ мест в кабине. Места продаются в $K$ тарифных классах, упорядоченных по убыванию цены: $p_1 > p_2 > \dots > p_K$. Период продаж задан шкалой дней до вылета (DFO) $d \in \{d_{\max}, \dots, 0\}$, при этом время идёт в сторону уменьшения $d$.

Спрос на класс $k$ за весь период продаж является случайной величиной $D_k$, его распределение характеризуется вектором параметров $\boldsymbol{\theta}_k$. Моменты поступления запросов задаются кривой накопления $G_k(d) \in [0,1]$: доля спроса класса $k$, поступившая к моменту $d$. Спрос разных классов считается независимым.

Политика управления задаётся вектором лимитов защиты $\mathbf{y} = (y_1, \dots, y_{K-1})$. Величина $y_k$ обозначает число мест, защищённых для классов $1, \dots, k$ от продажи в классах с более высокими номерами. Запрос в класс $j$ принимается, если число свободных мест превышает $y_{j-1}$, и отклоняется в противном случае.

Выручка рейса при реализации спроса $\omega$ и политике $\mathbf{y}$ равна
\begin{equation}
  R(\mathbf{y}, \omega) = \sum_{k=1}^{K} p_k \, x_k(\mathbf{y}, \omega),
  \label{eq:revenue}
\end{equation}
где $x_k(\mathbf{y}, \omega)$~--- число проданных мест класса $k$, причём $\sum_k x_k \le C$.

Требуется решить две задачи.
\begin{enumerate}[label=\arabic*)]
  \item По историческим данным о бронированиях оценить параметры спроса $\boldsymbol{\theta}_k$ и кривые $G_k(d)$ с учётом цензурирования.
  \item Найти лимиты, максимизирующие ожидаемую выручку:
  \begin{equation}
    \mathbf{y}^{*} = \arg\max_{\mathbf{y}} \; \mathbb{E}\, R(\mathbf{y}, \omega).
    \label{eq:opt}
  \end{equation}
\end{enumerate}
Качество решения оценивается сравнением расчётной выручки по оптимизированной политике $\mathbf{y}^{*}$ с выручкой, фактически полученной на тех же рейсах.

% ============================================================
\section{Описание методов}

\subsection{Исходные данные}

Используются обезличенные инвентарные данные авиакомпании по бронированиям международного регулярного направления. Данные детализированы по тарифным классам и кабинам и привязаны к шкале DFO. Для каждого рейса известны накопленное число бронирований по классам на каждый день шкалы, моменты закрытия классов, итоговая загрузка и итоговая выручка. Количество рейсов в выборке определяется доступной глубиной хранения данных в инвентарной системе.

\subsection{Восстановление спроса}

Фактические продажи по классу $k$ ограничены закрытием класса, поэтому наблюдаемое число бронирований $b_k$ меньше или равно спросу $D_k$. Для восстановления спроса используется кривая накопления бронирований. Рейсы разделяются на две группы: на рейсах первой группы класс $k$ оставался открытым до вылета, и наблюдаемые бронирования совпадают со спросом; на рейсах второй группы класс был закрыт на дне $d_c$.

По рейсам первой группы строится эмпирическая кривая накопления $\hat{G}_k(d)$, а для рейсов второй группы спрос оценивается как
\begin{equation}
  \hat{D}_k = \frac{b_k(d_c)}{\hat{G}_k(d_c)},
  \label{eq:unconstrain}
\end{equation}
где $b_k(d_c)$~--- число бронирований класса $k$ к моменту закрытия. По совокупности значений $\hat{D}_k$ оцениваются параметры распределения спроса: среднее $\mu_k$ и стандартное отклонение $\sigma_k$ в случае нормального закона либо интенсивность $\lambda_k$ в случае пуассоновского. Выбор закона определяется по критерию согласия на исторической выборке.

\subsection{Имитационная модель продаж}

Имитационная модель воспроизводит процесс продаж на одном рейсе методом Монте-Карло. Один прогон модели для заданных лимитов $\mathbf{y}$ выполняется в следующем порядке.
\begin{enumerate}[label=\arabic*)]
  \item Для каждого класса $k$ разыгрывается реализация спроса $D_k$ из оценённого распределения.
  \item Запросы класса $k$ распределяются по шкале DFO в соответствии с кривой $\hat{G}_k(d)$.
  \item Запросы всех классов обрабатываются в порядке убывания $d$: запрос принимается, если выполняется условие защиты, описанное в разделе~III, и отклоняется в противном случае.
  \item По числу проданных мест вычисляется выручка~\eqref{eq:revenue}.
\end{enumerate}

Прогон повторяется $N$ раз. Ожидаемая выручка оценивается выборочным средним
\begin{equation}
  \hat{R}(\mathbf{y}) = \frac{1}{N} \sum_{i=1}^{N} R(\mathbf{y}, \omega_i),
  \label{eq:mc}
\end{equation}
а доверительный интервал уровня $1-\alpha$ строится как
\begin{equation}
  \hat{R}(\mathbf{y}) \pm t_{1-\alpha/2,\,N-1} \frac{s}{\sqrt{N}},
  \label{eq:ci}
\end{equation}
где $s$~--- выборочное стандартное отклонение выручки. Для корректного сравнения разных политик используются одни и те же реализации спроса (метод общих случайных чисел), что уменьшает дисперсию оценки разности выручек.

\subsection{Поиск оптимальных лимитов}

Максимизация~\eqref{eq:opt} выполняется по сетке допустимых значений $\mathbf{y}$ с условием монотонности $y_1 \le y_2 \le \dots \le y_{K-1} \le C$. В качестве начального приближения используются лимиты, рассчитанные по правилу Литтлвуда и эвристике EMSR~\cite{littlewood1972, belobaba1987}, затем вокруг них выполняется покоординатный поиск, на каждом шаге которого значение $\hat{R}(\mathbf{y})$ вычисляется по~\eqref{eq:mc}. Поиск завершается, когда улучшение выручки становится статистически незначимым с учётом интервала~\eqref{eq:ci}.

\subsection{Валидация модели}

Для проверки модели используются те же исторические рейсы, по которым оценивался спрос. Для каждого рейса сравниваются три величины: фактическая выручка $R^{\text{fact}}$, расчётная выручка при воспроизведении фактически применённых лимитов $\hat{R}(\mathbf{y}^{\text{fact}})$ и расчётная выручка при оптимизированных лимитах $\hat{R}(\mathbf{y}^{*})$.

Адекватность имитационной модели определяется относительным отклонением
\begin{equation}
  \delta = \frac{\hat{R}(\mathbf{y}^{\text{fact}}) - R^{\text{fact}}}{R^{\text{fact}}},
  \label{eq:delta}
\end{equation}
а эффект от оптимизации~--- величиной $\hat{R}(\mathbf{y}^{*}) - \hat{R}(\mathbf{y}^{\text{fact}})$ с доверительным интервалом. Если фактическая выручка попадает в доверительный интервал для $\hat{R}(\mathbf{y}^{\text{fact}})$, модель считается воспроизводящей реальный процесс продаж, и полученная по ней оценка эффекта оптимизации может использоваться для практических выводов.

% ============================================================
\begin{thebibliography}{9}

\bibitem{fridman2025} Фридман~Г.~М., Баранов~В.~М. Имитационная модель процесса продаж авиабилетов для продуктов одного рейса // Вестник Воронежского государственного университета. Серия: Экономика и управление. 2025. №~2. С.~17--25. DOI: 10.17308/econ.2025.2/13065.

\bibitem{nosova} Носова~Е.~В. Моделирование процесса управления доходами авиакомпании на уровне рейсов // Известия Российского государственного педагогического университета им.~А.~И.~Герцена. Вып.~141. С.~70--75.

\bibitem{lavrenyuk} Лавренюк~О.~В. Описание компьютерного симулятора для системы управления доходами авиакомпании // Известия Российского государственного педагогического университета им.~А.~И.~Герцена. Вып.~141. С.~75--80.

\bibitem{mozgovaya2013} Мозговая~К.~А. Математическое моделирование сверхлимитных продаж в управлении доходами авиакомпании: автореф. дис. \dots{} канд. экон. наук: 08.00.13. СПб.: Санкт-Петербургский государственный экономический университет, 2013. 23~с. EDN: ZOYPCT.

\bibitem{littlewood1972} Littlewood~K. Forecasting and control of passenger bookings // AGIFORS Symposium Proceedings. 1972. Vol.~12. P.~95--117.

\bibitem{belobaba1987} Belobaba~P.~P. Air travel demand and airline seat inventory management: PhD thesis. Cambridge, MA: Massachusetts Institute of Technology, 1987.

\bibitem{talluri2004} Talluri~K.~T., van~Ryzin~G.~J. The Theory and Practice of Revenue Management. New York: Springer, 2004.

\end{thebibliography}

\end{document}